\section{Masked-data inference: brief primer}
\label{sec:background}

We summarize the masked-data identifiability framework
\citep{towell2026masked,heitjan1991ignorability}. Readers familiar
with this work or with \citet{towell2026scrnacoarsening,
towell2026spatialcoarsening} can skim to \cref{sec:translation}.

Consider $m$ independent components in series. Each component $j$ has a
parametric distribution $f_j(\cdot;\theta_j)$ for its lifetime $T_{ij}$
in system $i$. The system fails at $T_i = \min_j T_{ij}$ with cause
$K_i = \arg\min_j T_{ij}$. We observe $T_i$ but not $K_i$ directly;
instead we observe a \emph{candidate set} $C_i \subseteq \{1,\ldots,m\}$
thought to contain the failed component.

The joint likelihood factorizes as
\begin{equation}
  f(t, k, c; \theta)
  = h_k(t; \theta_k) \prod_{\ell} R_\ell(t; \theta_\ell)
    \cdot \Prob_\theta\{C_i = c \mid T_i = t, K_i = k\},
  \label{eq:bg-joint}
\end{equation}
where $h_\ell$ is the component-$\ell$ hazard and $R_\ell$ its survival
function. The third factor is the unknown masking mechanism.

Three coarsening conditions allow the masking factor to drop from the
score:

\begin{condition}[Support, C1]
\label{cond:c1}
$\Prob\{K_i \in C_i\} = 1$.
\end{condition}

\begin{condition}[Symmetry, C2]
\label{cond:c2}
$\Prob_\theta\{C_i=c \mid T_i=t, K_i=j\} = \Prob_\theta\{C_i=c \mid
T_i=t, K_i=k\}$ for all $j,k \in c$.
\end{condition}

\begin{condition}[Parameter independence, C3]
\label{cond:c3}
$\Prob_\theta\{C_i=c \mid T_i, K_i\}$ does not depend on $\theta$.
\end{condition}

Under all three, the likelihood collapses to
\begin{equation}
  L_i(\theta) \propto \left[\prod_\ell R_\ell(t_i;\theta_\ell)\right]
  \cdot \left[\sum_{j \in c_i} h_j(t_i; \theta_j)\right].
  \label{eq:bg-c1c2c3}
\end{equation}

\citet[Theorem 8]{towell2026masked} establish identifiability:

\begin{theorem}[Identifiability under C1--C2--C3]
\label{thm:bg-id}
Under \cref{cond:c1,cond:c2,cond:c3} and standard regularity, full
column rank $m$ of the augmented candidate-set matrix $\tilde{C} \in
\{0,1\}^{n \times m}$ is \emph{sufficient} for $\theta$ to be
identifiable. Separability is necessary but not in general
sufficient (a separating mechanism can leave $\tilde{C}$
rank-deficient, e.g. $m=4$ with candidate sets
$\{1,2\},\{3,4\},\{1,3\},\{2,4\}$).
\end{theorem}

\emph{Singleton} candidate sets ($|c_i|=1$) are critical: each isolates
a single component contribution in \eqref{eq:bg-c1c2c3} and adds an
identifying row to $\tilde{C}$.

The condition that does the conceptual work in this paper is C2.
\Cref{cond:c2} requires the candidate set to be generated
symmetrically with respect to which masked cause is true. When the
masking mechanism depends on the latent cause in a way not captured by
the modeled structure, C2 fails, the masking factor no longer drops
from the score, and cause-level parameters become biased. This is the
\emph{informative coarsening} regime. The sensitivity of estimates to
C2 violations is studied directly in \citet{towell2026mdrelax}.

The framework has been ported to single-cell RNA sequencing zero
inflation \citep{towell2026scrnacoarsening} and to spatial
transcriptomics deconvolution \citep{towell2026spatialcoarsening}. We
now port it to electronic phenotyping. Two features distinguish the
phenotyping case. First, the latent cause is binary: a patient either
has condition $Z$ or does not. Second, and more important, the
coarsening mechanism is structurally informative. The coding process
depends on disease severity, encounter frequency, and billing
incentives, so C2 is not an approximation that holds after modeling
technical covariates; it is the condition that fails.
